\begin{publications}
\noindent
\textbf{发表与接收论文}
\renewcommand{\labelenumi}{[\arabic{enumi}]}
\begin{enumerate}
	\item \textbf{Xilong Pei}, Haifan Yin, Li Tan, Lin Cao, Zhanpeng Li, Kai Wang, Kun Zhang, Emil Bj\"ornson.
	RIS-aided wireless communications: Prototyping, adaptive beamforming, and indoor/outdoor field trials.
	IEEE Transactions on Communications, 2021, 69(12): 8627--8640. 
	（SCI源刊；IF：8.3；署名单位：华中科技大学；ESI高被引论文；第一作者）
	
	\item \textbf{Xilong Pei}, Haifan Yin, Li Tan, Lin Cao.
	A two-stage spatial-oversampling codebook and field trials of RIS-aided wireless communications.
	in: IEEE 35th International Symposium on Personal, Indoor and Mobile Radio Communications (PIMRC).
	Valencia, Spain, 2--5 Sep., IEEE, 2024: 1--6.
	（EI检索；署名单位：华中科技大学；第一作者）
	
	\item \textbf{Xilong Pei}, Haifan Yin, Lin Cao, Rongguang Song.
	Dynamic metasurface antennas with discrete phase shifts: Performance analysis and beamforming methods.
	IEEE Transactions on Communications, 2025, 73(12): 15146--15159.
	（SCI源刊；IF：8.3；署名单位：华中科技大学；第一作者）
	
	\item \textbf{Xilong Pei}, Haifan Yin, Lin Cao, Rongguang Song, Yuanwei Liu.
	Geometry-guided pinching antenna positioning for enhanced localization accuracy.
	IEEE Wireless Communications Letters, 2026, 15: 2749--2753.
	（SCI源刊；IF：5.5；署名单位：华中科技大学；第一作者）
	
	\item Lin Cao, Haifan Yin, \textbf{Xilong Pei}, Rongguang Song.
	Initial excitation estimation for DMAs via amplitude-only measurements.
	IEEE Transactions on Antennas and Propagation, 2026, \textcolor{red}{xxxx--xxxx}.
	（SCI源刊；IF：5.8；署名单位：华中科技大学；学生第二作者）
	
	\item Lin Cao, Haifan Yin, \textbf{Xilong Pei}, Rongguang Song, Yuanwei Liu.
	Multi-pin and movable-pin enabled discretely-activated pinching antenna systems.
	IEEE Wireless Communications Letters, 2026, 15: 590--594. 
	（SCI源刊；IF：5.5；署名单位：华中科技大学；学生第二作者）
	
	\item Rongguang Song, Haifan Yin, \textbf{Xilong Pei}, Lin Cao, Taorui Yang, Xue Ren, Yuanwei Liu.
	Design and prototyping of filtering active STAR-RIS with adjustable power splitting.
	IEEE Transactions on Antennas and Propagation, 2025, 73(9): 6462--6476.
	（SCI源刊；IF：5.8；署名单位：华中科技大学；学生第二作者）
	
	\item Lin Cao, Haifan Yin, Li Tan, \textbf{Xilong Pei}.
	RIS with insufficient phase shifting capability: Modeling, beamforming, and experimental validations.
	IEEE Transactions on Communications, 2024, 72(9): 5911--5923.
	（SCI源刊；IF：8.3；署名单位：华中科技大学；学生第二作者）
	
	\item Jiangfeng Hu, Haifan Yin, Li Tan, Lin Cao, \textbf{Xilong Pei}.
	RIS-aided wireless communications: Can RIS beat flat metal plate?.
	IEEE Transactions on Vehicular Technology, 2024, 73(10): 15704--15708.
	（SCI源刊；IF：7.1；署名单位：华中科技大学）
	
	\item Lin Cao, Haifan Yin, \textbf{Xilong Pei}, Li Tan.
	RIS with practical reflection coefficients: Modeling and experimental measurements.
	in: 18th European Conference on Antennas and Propagation (EuCAP). 
	Glasgow, United Kingdom, 17--22 Mar., IEEE, 2024: 1--5.
	（EI检索；署名单位：华中科技大学；学生第二作者）
	
	\item Zipeng Wang, Li Tan, Haifan Yin, Kai Wang, \textbf{Xilong Pei}, David Gesbert.
	A received power model for reconfigurable intelligent surface and measurement-based validations.
	in: IEEE 22nd International Workshop on Signal Processing Advances in Wireless Communications (SPAWC).
	Lucca, Italy, 27--30 Sep., IEEE, 2021: 561--565.
	（EI检索；署名单位：华中科技大学）
\end{enumerate}

\noindent
\textbf{专\hspace{2em}利}
\renewcommand{\labelenumi}{[\arabic{enumi}]}
\begin{enumerate}
	\item \textbf{裴熙隆}, 尹海帆, 谭力, 张昆, 王锴, 陆骋.
	一种应用于智能超表面的控制方法及智能超表面的控制器.
	中国, 发明专利, CN202110380307.9.
	（署名单位：华中科技大学；学生第一发明人）
	
	\item 尹海帆, \textbf{裴熙隆}, 曹琳, 宋镕光.
	用于动态超表面天线的离散相位最优波束成形方法及系统.
	中国, 发明专利, CN202511446755.9.
	（署名单位：华中科技大学；学生第一发明人）
	
	\item 尹海帆, 曹琳, \textbf{裴熙隆}, 谭力.
	一种基于幅度测量的幅度-相位耦合型天线阵列初始激励估计方法.
	中国, 发明专利, CN202610055887.7.
	（署名单位：华中科技大学；学生第二发明人）
	
	\item 尹海帆, 宋镕光, 吴集旺, \textbf{裴熙隆}.
	有源透射型智能超表面及应用于此智能超表面的配套电路.
	中国, 发明专利, CN202411350355.3.
	（署名单位：华中科技大学；学生第三发明人）
	
	\item 曹琳, 王锴, 尹海帆, 陆骋, \textbf{裴熙隆}, 张昆.
	毫米波智能超表面单元及毫米波智能超表面.
	中国, 发明专利, CN202110380371.7.
	（署名单位：华中科技大学；学生第四发明人）
\end{enumerate}

\end{publications}
